\documentclass[10pt,a4paper]{amsart}
\usepackage{amsmath,amssymb,amsthm,graphicx}
% Retained minimal preamble for the Stacks Project chapter more-algebra.tex (More on Algebra).
% The chapter file pinned-original-source.tex carries only \input{preamble} on its first line, so no
% preamble text lives inside the chapter itself.  The Stacks Project uses one single shared preamble for
% every chapter; the snapshot stages it identically as sources/stacks-chow/preamble.tex and
% sources/stacks-etale-cohomology/preamble.tex.  The block below is copied mechanically, line for line,
% from sources/stacks-chow/preamble.tex lines 14-18 (the Xy-pic setup) and 154-217 (the native theorem
% declarations together with the mathematical macros).  Nothing was added, removed or redefined; the
% wrapper class and the amsmath/amssymb/amsthm/graphicx base packages are supplied by the assembler.
% Chapter: more-algebra.tex -- More on Algebra.
% For commutative diagrams we use Xy-pic
\usepackage[all]{xy}

% We use 2cell for 2-commutative diagrams.
\xyoption{2cell}

% Theorem environments.
%
\theoremstyle{plain}
\newtheorem{theorem}[subsection]{Theorem}
\newtheorem{proposition}[subsection]{Proposition}
\newtheorem{lemma}[subsection]{Lemma}

\theoremstyle{definition}
\newtheorem{definition}[subsection]{Definition}
\newtheorem{example}[subsection]{Example}
\newtheorem{exercise}[subsection]{Exercise}
\newtheorem{situation}[subsection]{Situation}

\theoremstyle{remark}
\newtheorem{remark}[subsection]{Remark}
\newtheorem{remarks}[subsection]{Remarks}

\numberwithin{equation}{subsection}

% Macros
%
\def\lim{\mathop{\mathrm{lim}}\nolimits}
\def\colim{\mathop{\mathrm{colim}}\nolimits}
\def\Spec{\mathop{\mathrm{Spec}}}
\def\Hom{\mathop{\mathrm{Hom}}\nolimits}
\def\Ext{\mathop{\mathrm{Ext}}\nolimits}
\def\SheafHom{\mathop{\mathcal{H}\!\mathit{om}}\nolimits}
\def\SheafExt{\mathop{\mathcal{E}\!\mathit{xt}}\nolimits}
\def\Sch{\mathit{Sch}}
\def\Mor{\mathop{\mathrm{Mor}}\nolimits}
\def\Ob{\mathop{\mathrm{Ob}}\nolimits}
\def\Sh{\mathop{\mathit{Sh}}\nolimits}
\def\NL{\mathop{N\!L}\nolimits}
\def\CH{\mathop{\mathrm{CH}}\nolimits}
\def\proetale{{pro\text{-}\acute{e}tale}}
\def\etale{{\acute{e}tale}}
\def\QCoh{\mathit{QCoh}}
\def\Ker{\mathop{\mathrm{Ker}}}
\def\Im{\mathop{\mathrm{Im}}}
\def\Coker{\mathop{\mathrm{Coker}}}
\def\Coim{\mathop{\mathrm{Coim}}}

% Boxtimes
%
\DeclareMathSymbol{\boxtimes}{\mathbin}{AMSa}{"02}

%
% Macros for moduli stacks/spaces
%
\def\QCohstack{\mathcal{QC}\!\mathit{oh}}
\def\Cohstack{\mathcal{C}\!\mathit{oh}}
\def\Spacesstack{\mathcal{S}\!\mathit{paces}}
\def\Quotfunctor{\mathrm{Quot}}
\def\Hilbfunctor{\mathrm{Hilb}}
\def\Curvesstack{\mathcal{C}\!\mathit{urves}}
\def\Polarizedstack{\mathcal{P}\!\mathit{olarized}}
\def\Complexesstack{\mathcal{C}\!\mathit{omplexes}}
% \Pic is the operator that assigns to X its picard group, usage \Pic(X)
% \Picardstack_{X/B} denotes the Picard stack of X over B
% \Picardfunctor_{X/B} denotes the Picard functor of X over B
\def\Pic{\mathop{\mathrm{Pic}}\nolimits}
\def\Picardstack{\mathcal{P}\!\mathit{ic}}
\def\Picardfunctor{\mathrm{Pic}}
\def\Deformationcategory{\mathcal{D}\!\mathit{ef}}

% --- appended 2026-10-12 (Chief): the project own definitions of the reference/slogan/history
% display environments, copied mechanically from the shared Stacks preamble lines 8-12
% (\usepackage{verbatim} plus the three \newenvironment lines, which the chapter preamble keeps
% verbatim from the project preamble). Nothing added, removed or redefined: these environments
% carry no mathematical content and the project itself typesets them as comments. Without these
% three lines a counted statement whose own text uses them fails with Environment slogan undefined.
\usepackage{verbatim}
\newenvironment{reference}{\comment}{\endcomment}
\newenvironment{slogan}{\comment}{\endcomment}
\newenvironment{history}{\comment}{\endcomment}
\usepackage{cleveref}
\crefname{originalequation}{original equation}{original equations}
\Crefname{originalequation}{Original equation}{Original equations}
\makeatletter
\newlabel{equation-BL-cech-re}{{0F1Q}{source}{}{sourcepointer-1}{}}
\newlabel{equation-BL-cech-re@cref}{{[originalequation][0][]0F1Q}{[1][0][]source}{}{}{}}
\makeatother
\crefname{originallemma}{original lemma}{original lemmas}
\Crefname{originallemma}{Original lemma}{Original lemmas}
\makeatletter
\newlabel{lemma-same-quotients}{{0BNJ}{source}{}{sourcepointer-2}{}}
\newlabel{lemma-same-quotients@cref}{{[originallemma][0][]0BNJ}{[1][0][]source}{}{}{}}
\makeatother
\title{Exactness of the Cech sequence of the ring itself is equivalent to bijectivity on f-power torsion, giving a glueing pair criterion}
\author{The Stacks Project Authors}
\date{stacks-project:more-algebra.tex:local-snapshot-20261012; GNU FDL 1.2 or later}
\begin{document}
\maketitle
\setcounter{subsection}{1}
\noindent\textit{Editorial note.} All delivered mathematical text is the author's own, copied from the pinned Stacks Project chapter file more-algebra.tex; no statement or proof step was rewritten, reordered or removed. The counted statement is the labeled lemma environment at source lines 25254-25270 (canonical tag 0BNR) and the counted proof is the proof environment at source lines 25272-25298, which immediately follows it and is its own proof.\par\medskip
\noindent\textit{Editorial note.} The statement refers to the displayed Cech sequence by its label equation-BL-cech-re, which lies outside the counted spans; that label is carried as a navigation stub so the reference resolves in the compiled document.\par\medskip
\noindent\textit{Editorial note.} The preamble span runs to line 90, the current end of the retained preamble, which includes the appended environment definitions; the counted text does not use them, and the span is taken whole so that the file version is unambiguous.\par\medskip
\noindent\textit{Editorial note.} The source URL is the source-level Stacks Project URL, as required when more than one card is produced from one source; the per-claim tag is carried by claim\_key and by this note. The counted statement and proof cite no bibliography entry, so this card needs no bibliography.\par\medskip
\noindent\textit{Allowed standard prerequisites (editorial).}\begin{itemize}
\item Every element of the localization of the target ring at the element can be written with a denominator, and the difference between a lift and the numerator can be absorbed by a multiple of a power of the element using that the quotients by the powers of the element agree. used at proof lines 25273-25276. Gives exactness on the right by producing an explicit pair mapping to a given element of the localization.
\item The kernel of the localization map of a ring is the ideal of elements annihilated by a power of the element. used at proof line 25279. Gives the left exactness criterion.
\item The last statement of the lemma is a special case of the penultimate one, with the completion of the ring in the role of the second ring. used at proof lines 25296-25297. Transfers the glueing pair criterion from an arbitrary second ring to the completion.
\end{itemize}\medskip
\ifdefined\hypertarget\hypertarget{sourcepointer-1}{}\fi
\noindent\textit{Source reference.} Original equation 0F1Q: sources/stacks-more-algebra/pinned-original-source.tex, source lines 25238--25238. The reference points to the displayed Cech sequence that the statement is about; it names the object under discussion and the element computations deciding its exactness are inside the counted span.\par\medskip
\ifdefined\hypertarget\hypertarget{sourcepointer-2}{}\fi
\noindent\textit{Source reference.} Original lemma 0BNJ: sources/stacks-more-algebra/pinned-original-source.tex, source lines 25129--25129. The reference is invoked at proof lines 25296-25297 to obtain the completion case as a special case of the penultimate statement; it is the final step and the rest of the proof is inside the counted span.\par\medskip
\begin{lemma}
\label{lemma-same-f-torsion}
Let $R$ be a ring, let $f \in R$, and let $R \to R'$ be a ring map
which induces isomorphisms $R/f^nR \to R'/f^nR'$ for $n > 0$.
The sequence (\ref{equation-BL-cech-re}) is
\begin{enumerate}
\item exact on the right,
\item exact on the left if and only if $R[f^\infty] \to R'[f^\infty]$
is injective, and
\item exact in the middle if and only if $R[f^\infty] \to R'[f^\infty]$
is surjective.
\end{enumerate}
In particular, $(R \to R', f)$ is a glueing pair if and only if
$R[f^\infty] \to R'[f^\infty]$ is bijective. For example, $(R, f)$
is a glueing pair if and only if
$R[f^\infty] \to R^\wedge[f^\infty]$ is bijective.
\end{lemma}

\begin{proof}
Let $x \in R'_f$. Write $x = x'/f^n$ with $x' \in R'$.
Write $x' = x'' + f^n y$ with $x'' \in R$ and $y \in R'$.
Then we see that $(y, -x''/f^n)$ maps to $x$. Thus (1) holds.

\medskip\noindent
Part (2) follows from the fact that $\Ker(R \to R_f) = R[f^\infty]$.

\medskip\noindent
If the sequence is exact in the middle, then elements of the form
$(x, 0)$ with $x \in R'[f^\infty]$ are in the image of
the first arrow. This implies that $R[f^\infty] \to R'[f^\infty]$
is surjective. Conversely, assume that $R[f^\infty] \to R'[f^\infty]$
is surjective. Let $(x, y)$ be an element in the middle
which maps to zero on the right. Write $y = y'/f^n$ for some $y' \in R$.
Then we see that $f^n x - y'$ is annihilated by some power of $f$ in
$R'$. By assumption we can write $f^nx - y' = z$ for some
$z \in R[f^\infty]$. Then $y = y''/f^n$ where $y'' = y' + z$
is in the kernel of $R \to R/f^nR$. Hence we see that $y$ can be
represented as $y'''/1$ for some $y''' \in R$. Then
$x - y'''$ is in $R'[f^\infty]$. Thus $x - y''' = z' \in R[f^\infty]$.
Then $(x, y'''/1) = (y''' + z', (y''' + z')/1)$ as desired.

\medskip\noindent
The last statement of the lemma is a special case of the penultimate statement
by Lemma \ref{lemma-same-quotients}.
\end{proof}

\end{document}
